% ============================================================================
% 第 1 章 引言与问题背景
% ============================================================================
\section{引言与问题背景}
\label{sec:introduction}

\subsection{工业背景与研究意义}

随着智能制造、电商履约与柔性供应链的快速发展，传统人工拣选仓储模式在吞吐效率、空间利用率与人力成本等方面逐渐逼近瓶颈，自动化立体仓库（Automated Storage and Retrieval System, AS/RS）及多穿梭无人车（Automated Guided Vehicle, AGV）集群已成为现代智慧物流的核心基础设施。在典型的零件与货品立体仓库中，各类物料按批次和库位分布于密集排列的多层货架内；当上游生产系统或订单管理系统持续下发取货、出库与补货请求时，调度中心需要实时指挥多辆自主运行的 AGV 在有限通道空间内完成“空驶就位—货架取货—送达交接站—待命或最终返库”的闭环作业。

相比单车路径规划或离线静态排程，多取货任务并发下的仓储 AGV 调度在理论与工程上均属于典型的 NP-hard 组合优化与多智能体协同决策问题~\cite{ma2017lifelong, stern2019mapf}。在密集货架与狭窄巷道布局中，简单套用“最近车辆就近派单”或“静态最短路独立寻路”往往会引发严重的连锁退化：多辆车在同一交接站出口或主通道汇流形成拥堵与死锁，局部最优的即时派单破坏后续波次的任务均衡，进而拉长整体完工时间（Makespan）与订单响应延迟。因此，如何在严格满足物理避碰与仓储工艺规则的前提下，实现高吞吐、低延迟、低能耗且具备良好计算实时性的多 AGV 智能调度，具有直接的工程价值与重要的算法研究意义。

\subsection{多 AGV 并发调度的核心技术难点}

结合现代立体仓库的实际运行机理，多任务并发下的 AGV 智能调度主要面临以下四个相互交织的核心难点：

\begin{enumerate}
  \item \textbf{任务指派与时空路径规划的强耦合性}。车辆执行某项取货任务的真实代价并非静态几何距离，而是取决于当前在途车队的时空占位分布与巷道拥堵状态。若将任务分配与路径规划完全割裂，上层基于曼哈顿距离或静态最短路估算的“最优指派”在底层路径展开时可能遭遇大量等待与绕行；反之，若将全部车辆的离散时空状态联合纳入单一全局求解器，搜索空间随车辆数、订单数与规划时域呈指数级膨胀，难以满足在线滚动决策的时间约束~\cite{li2021rhcr}。
  \item \textbf{多层立体货架与复杂作业规则的硬约束}。真实零件仓库并非理想二维空旷栅格，而是包含单面多层货架、背靠背混合货架、单向/共享窄巷道、立柱障碍以及容量受限的交付工作站。任务不仅具有位置属性，还携带货架层高（直接决定升降装货耗时）及同面货架的层序解锁关系——同一货架面中必须遵循既定取货顺序，且交付站卸货后须及时驶离以释放终端资源。此外，入库与出库流可能部分重叠，空闲车辆在波次间隙需合理待命，并在订单流显式关闭后安全返库。
  \item \textbf{在线动态到达与因果信息边界}。在持续运行的真实订单流中，调度器无法预知未来尚未发布的订单，必须在严格遵守因果信息边界的前提下开展在线滚动决策。高负载高峰期的瞬时积压与低负载平峰期的零散到达交替出现，要求调度策略在“即时响应”与“适度协同筹划”之间取得动态平衡，同时避免因过度拟合特定历史波次而在未知订单窗口上产生泛化退化。
  \item \textbf{多目标权衡与物理可行性保障}。仓储运维通常同时关注三类相互制约的性能指标：决定系统整体周期效率的\emph{最终完工时间}（$T$）、反映单笔订单服务质量的\emph{平均响应耗时}（$R$），以及关联设备磨损与能耗的\emph{总行驶里程}（$D$）。此外，任何优化算法产生的调度方案都必须保证 100\% 满足顶点无碰撞、边交换无冲突、连续装卸时长与终端资源独占等底层物理约束，否则在物理系统中无法落地执行。
\end{enumerate}

\subsection{研究目标与本文主要贡献}

针对上述挑战，本项目以“理论建模清晰、算法演进有据、工程实现可复现”为原则，构建了面向多取货任务并发的仓库 AGV 智能调度与仿真验证体系。围绕课题要求，本文的主要研究工作与贡献可概括如下：

\begin{itemize}
  \item \textbf{构建分层解耦且严格可验的调度仿真架构}：设计并实现了纯 Python 标准库（零第三方重型依赖）的离散时空仿真执行内核与独立物理验解器，打通了“多工况建模—在线调度—冲突校验—可视化回放”的全链路，并支持交互式 Web 动态仿真回放与一键复现。
  \item \textbf{形成“通用多工况引擎”与“真实工业订单流”双主线协同研究}：一方面在自研通用引擎（\code{src/agv/}）中系统支持 Legacy 规则栅格、V4 分层单面货架及 V9/V9.1 混合背靠背货架布局，提出基于安全时间区间路径规划（SIPP）、拥堵关联大邻域修复（LNS）与高积压影子试派（BCRD）的模块化算法族；另一方面接入 SLAPStack/WEPAStacks 公开真实入出库订单流~\cite{pfrommer2022slapstack}，在 $62\times 47$ 大规模立体仓库网格（6 车、5 层、900\,s 订单窗口）下开展了多轮对照与跨窗口稳健性研究。
  \item \textbf{收敛出高效稳健的最终调度方案 \code{gate}}：在统一固定尺度线性目标函数下，通过问题驱动的多轮消融与跨窗口复评，最终定型“事件门控滚动派单＋软分区联合遗传算法（GA）＋预约表 SIPP 物理预演”方案（\code{gate}）。在涵盖 561 次独立验解运行的系统实验中，\code{gate} 在全部既有验证窗口上稳定优于 \code{fcfs}、\code{eta} 与 \code{hungarian} 基线，并通过事件门控大幅压缩低负载场景的无效搜索触发。
\end{itemize}

\subsection{本文结构安排}

本文其余部分组织如下：第~\ref{sec:methodology} 章给出系统总体架构、三层协同机制与最终采用方案概要；第~\ref{sec:evolution} 章沿问题驱动的视角系统梳理两条主线的思路推进、算法演化与关键消融决策；第~\ref{sec:modeling} 章阐述仓库几何、物理约束与评价指标的形式化数学模型；第~\ref{sec:algorithms} 章详述核心调度与路径规划算法设计；第~\ref{sec:experiments} 章展示跨场景、多随机种子及超参数敏感性实验结果与对比分析；第~\ref{sec:conclusion} 章总结全文工作并展望后续拓展方向。
